\subsection*{章末の案内語}

この表は、第5章で説明した語のうち、名前から宛先までパケットを運ぶ主経路に必要な語だけを示します。

\noindent\begin{tabularx}{\linewidth}{@{}>{\bfseries\raggedright\arraybackslash\hyphenpenalty=10000}p{32mm}>{\raggedright\arraybackslash}p{35mm}Y@{}}
  \toprule
  この章の語 & 最初に説明した節 & 一言で戻る \\
  \midrule
  \guideterm{パケット} & 回線交換とパケット交換 & データを分けて網の中を運ぶ単位です。 \\
  \guideterm{IP} & IPアドレスとCIDR & パケットを宛先アドレスへ届ける共通層です。 \\
  \guideterm{IPアドレス} & IPアドレスとCIDR & インターネット層で通信端点を識別する番号です。 \\
  \guideterm{CIDR} & IPアドレスとCIDR & アドレスの先頭から何ビットを網として扱うか示します。 \\
  \guideterm{ルーティング} & ルーティング & 宛先に応じて次の転送先を選びます。 \\
  \guideterm{AS／BGP} & 網の網／経路制御の自律性 & 組織単位の網と、その間で経路情報を交換する仕組みです。 \\
  \guideterm{DNS} & DNS & ドメイン名から通信に使う情報を調べます。 \\
  \guideterm{リゾルバ} & DNS & 利用者に代わって名前解決を進めます。 \\
  \guideterm{TTL} & DNS & DNS応答を通常のキャッシュとして保持する時間を示します。 \\
  \bottomrule
\end{tabularx}
